% Separate continuation by /root/frobenius; C.4x.
\paragraph{C.4x (Proved: full real multiplication in arbitrary positive characteristic).}
Let \(k\) be algebraically closed of characteristic \(p>0\), and
let \(A/k\) have dimension \(g\). Suppose that the full integer
ring of a totally real field of degree \(g\) acts:
\[
                   \mathcal O_F\hookrightarrow\operatorname{End}(A),
                 \qquad [F:\mathbb Q]=g.
\]
We assume that \(A\) has a polarization whose Rosati involution
fixes \(F\) pointwise, equivalently an \(\mathcal O_F\)-linear
polarization. Then the integer \(\delta(A)\) of C.4j is a power
of \(p\), with exponent possibly zero. If, in addition,
\begin{equation}\label{eq:c4-hb-lie-condition}
 [\operatorname{Lie}(A)]
       =[\mathcal O_F\otimes_{\mathbb Z} k]
 \quad\text{in }G_0(\mathcal O_F\otimes_{\mathbb Z}k),
\end{equation}
where \(G_0\) is the Grothendieck group of finite modules, then
\[
                              \delta(A)=1.
\]
Condition \eqref{eq:c4-hb-lie-condition} is automatic if \(p\)
is totally ramified in \(F\).
There is no ordinarity or supersingularity hypothesis, and
\(p=2\) and \(p=3\) are included.

Consequently, in dimensions two and three respectively,
\begin{equation}\label{eq:c4-hb-splitting-bound}
 s_J(A)\mid p^{v_p(2)}
       \quad(g=2),\qquad
 s_J(A)\mid p^{v_p(6)}
       \quad(g=3).
\end{equation}
Every such surface in characteristic different from two, and
every such threefold in characteristic at least five, is an
actual direct summand of a product of Jacobians.
In every characteristic, including the remaining two small
primes, \eqref{eq:c4-hb-lie-condition} gives the same affirmative
conclusion for \(g\le3\). In particular this holds when the
characteristic is totally ramified in the acting real field.

\begin{proof}
We first record the exact input concerning integral
polarizations. In the original published paper,
\citet[Lemma 2.7, pp.~2117--2118]{frC4Yu2003HB}
proves that, for each rational prime \(\ell\ne p\), an abelian
\(\mathcal O_F\)-variety of dimension \([F:\mathbb Q]\) has an
\(\mathcal O_F\)-linear polarization of degree prime to \(\ell\).
His statement even allows an Artinian base; we use only the
algebraically closed field case. We explain the local
calculation, including the integrality step.

Let \(\lambda_0\) be the given compatible polarization. Write
\[
 \mathcal O_F\otimes\mathbb Z_\ell=\prod_{w\mid\ell}\mathcal O_w,
       \qquad T_\ell(A)=\bigoplus_{w\mid\ell}T_w .
\]
Each \(T_w\) is free of rank two over \(\mathcal O_w\).
The pairing of \(\lambda_0\) on \(T_w\) is the trace of an
alternating \(\mathcal O_w\)-bilinear pairing with values in
the local codifferent \(\mathfrak D_w^{-1}\).
After choosing a basis of \(T_w\), let \(c_w\ge0\) be the
valuation of its coefficient relative to a generator of
\(\mathfrak D_w^{-1}\). Thus this trace pairing is perfect
precisely when \(c_w=0\).

Choose \(a\in\mathcal O_F[1/\ell]\) with
\(\operatorname{ord}_w(a)=-c_w\) for every \(w\mid\ell\),
and with \(a\) totally positive. Such a choice follows
by prescribing the finitely many local valuations and
then adding a sufficiently large positive integer divisible
by a sufficiently high power of \(\ell\): the valuations
are unchanged, and all real embeddings become positive.
The resulting rational homomorphism
\(\lambda_0a:A\dashrightarrow A^\vee\) has no denominators
away from \(\ell\), and its \(\ell\)-adic pairing is perfect.
To check that it is an actual homomorphism, write
\(a=b/\ell^n\) with \(b\in\mathcal O_F\).
Integrality of its \(\ell\)-adic map says that \(\lambda_0b\)
annihilates the finite group scheme \(A[\ell^n]\), which is
étale because \(\ell\ne p\). Hence \(\lambda_0b\) factors
through \([\ell^n]\), as required. Total positivity makes
this symmetric homomorphism a polarization, exactly as
in the proof of Yu's lemma. Its perfect Tate pairing
makes its degree prime to \(\ell\).

Choose an ample line bundle \(L_\ell\) inducing this
polarization. The identity
\[
                   \deg\varphi_{L_\ell}=\chi(L_\ell)^2
\]
shows that \(\ell\nmid\chi(L_\ell)\).
We obtain this separately for every \(\ell\ne p\);
therefore the positive gcd \(\delta(A)\) has no prime
divisor other than \(p\).

At the characteristic prime we use
\citet[Proposition 2.8, p.~2118]{frC4Yu2003HB}.
It identifies \eqref{eq:c4-hb-lie-condition} with existence
of a compatible polarization of degree prime to \(p\)
(called a separable polarization there), and also with
the Deligne--Pappas and determinant conditions.
The notation in Sections 2--6 of that paper initially
assumes one prime of \(F\) over \(p\); the reduction to
the separate local components is explicitly stated on
p.~2109 and in Section 2.1, p.~2114.
We use its intrinsic condition (3), namely
\eqref{eq:c4-hb-lie-condition}, so the assertion applies
without that simplifying assumption.
The resulting ample \(L_p\) has \(p\nmid\chi(L_p)\).
Together with the preceding prime-by-prime construction,
this gives \(\delta(A)=1\).

For clarity, total ramification supplies the intrinsic
Lie condition without a smoothness assumption on a
Hilbert moduli space. If \(p\mathcal O_F=\mathfrak p^g\),
then \(R=\mathcal O_F\otimes_{\mathbb Z}k\) is an
Artinian local \(k\)-algebra of length \(g\), with residue
field \(k\). Its only simple finite module is \(k\).
Consequently \(G_0(R)\simeq\mathbb Z\), with the class
of a finite module equal to its \(k\)-dimension.
Both \(R\) and \(\operatorname{Lie}(A)\) have dimension
\(g\); hence their classes agree. This proves
\eqref{eq:c4-hb-lie-condition} also in wild
characteristic and for a nonfree Lie module.

Finally, C.4j gives \(s_J(A)\mid\delta(A)\) in dimensions
at most three. Combining this with the improved
universal bounds \(s_J(A)\mid2\) for surfaces and
\(s_J(A)\mid6\) for threefolds proves
\eqref{eq:c4-hb-splitting-bound}.
Whenever the resulting gcd is one, the finite
factorizations in C.4j and the universal factorization
combine by Bézout to give one actual retraction through
a product of Jacobians. Thus the conclusion is integral.
\end{proof}

The Lie condition is used only as a sufficient condition
for \(\delta(A)=1\); its failure is not asserted to
obstruct a Jacobian-product retraction. Nor does this
argument supply a real multiplication action on an
arbitrary supersingular threefold. The generic
\(\operatorname{End}(A)=\mathbb Z\), type-\((1,2)\)
surface and arbitrary supersingular threefolds in
characteristics two and three therefore retain the
previously isolated integral-cycle question.
